%% ---------------------------------------------------------------------------
%% Two in Three Hundred -- package-hallucination exposure, and the extractor
%% Journal of Information Security and Applications (JISA) -- Elsevier preprint
%%
%% Build:  python3 build.py && pdflatex main && bibtex main && pdflatex main x2
%%
%% Every number in the tables and macros is generated by make_tables.py from
%% ../results/. Do not edit main.tex or tables/*.tex by hand.
%% ---------------------------------------------------------------------------
\documentclass[preprint,12pt]{elsarticle}

\usepackage{booktabs}
\usepackage{array}

%% Ngat dong: ten file va dinh danh dai trong \texttt la mot token lien,
%% khong co cho ngat nen bi day ra ngoai le cot.
\sloppy
\emergencystretch=3em
\tolerance=2000

\usepackage{graphicx}
\usepackage{amsmath}
\usepackage{amssymb}
\usepackage{url}
\usepackage[htt]{hyphenat}
%% OT1 thay vi T1: Computer Modern ma hoa T1 (ec) chi co ban bitmap
%% Metafont tren he nay, sinh ra chu Type 3 nhoe trong PDF. OT1 da co
%% san ban Type 1 vector. T5 (tieng Viet) dung VNR, cung la Type 1.
\usepackage[T5,OT1]{fontenc}

\newcommand{\vn}[1]{{\fontencoding{T5}\selectfont #1}}

\newcommand{\pendingfig}[2]{%
  \IfFileExists{figures/#1.pdf}%
    {\includegraphics[width=\linewidth]{figures/#1.pdf}}%
    {\fbox{\parbox[c][#2][c]{\dimexpr\linewidth-2\fboxsep-2\fboxrule\relax}%
       {\centering\sffamily\small
        [\,\texttt{\detokenize{figures/#1.pdf}}\,]\par\vspace{3pt}
        \footnotesize hand-drawn diagram --- drop the PDF in
        \texttt{figures/} under this exact name\par}}}}

\journal{IEEE Transactions on Software Engineering}

%%MACROS%%

\usepackage{orcidlink}
\begin{document}

\begin{frontmatter}

\title{Two in Three Hundred: Measuring Package-Hallucination Exposure When One
Invisible Character Is Worth Half the Signal}

\author[rmit]{Son X. Ha\,\orcidlink{0000-0002-5336-0566}}
\ead{ha.son@rmit.edu.vn}
\author[hcmut]{Phat T. Tran-Truong\,\orcidlink{0000-0003-3199-6333}\fnref{js}}
\ead{phatttt@hcmut.edu.vn}
\author[hcmut]{Xuan-Bach Le\,\orcidlink{0009-0003-6848-5403}\fnref{js}}
\ead{lexuanbach@hcmut.edu.vn}
\author[fptct]{Trung Phan Hoang Tuan\,\orcidlink{0009-0007-2430-8199}}
\ead{trungpht@fe.edu.vn}
\author[fpt]{Tung Q. Nguyen\,\orcidlink{0009-0003-5046-5535}}
\ead{NguyenQuangTung0902@gmail.com}
\author[fptct]{Tran Gia Huy\,\orcidlink{0009-0002-3085-1255}}
\ead{trangiahuyh26@gmail.com}
\author[fptct]{Nhan Nguyen Tran Kha\,\orcidlink{0009-0004-2709-2061}}
\ead{NhanNTKCE200304@gmail.com}
\author[fptct]{Thuc Nguyen Minh\,\orcidlink{0009-0000-4810-3613}}
\ead{ThucNMCE200458@gmail.com}
\author[tdtu1,tdtu2]{Nguyen Ngoc Phien\,\orcidlink{0000-0001-9152-6208}\corref{cor}}
\ead{nguyenngocphien@tdtu.edu.vn}

\cortext[cor]{Corresponding author.}
\fntext[js]{P. T. Tran-Truong and X.-B. Le contributed equally to this work and are joint second authors.}

%% Keep affiliation values plain ASCII on one line: elsarticle parses them as
%% key/value options and a font-switching macro inside a value breaks it.
\affiliation[hcmut]{organization={Faculty of Computer Science and Engineering, Ho Chi Minh City University of Technology (HCMUT), VNU-HCM}, city={Ho Chi Minh City}, country={Vietnam}}
\affiliation[rmit]{organization={The Business School, RMIT University}, city={Ho Chi Minh City}, country={Vietnam}}
\affiliation[fpt]{organization={FPT University, Ho Chi Minh City Campus}, city={Ho Chi Minh City}, country={Vietnam}}
\affiliation[fptct]{organization={FPT University, Can Tho Campus}, city={Can Tho}, country={Vietnam}}
\affiliation[tdtu1]{organization={Center for Applied Information Technology \& Faculty of Information Technology, Ton Duc Thang University}, city={Ho Chi Minh City}, country={Vietnam}}
\affiliation[tdtu2]{organization={Faculty of Information Technology, Ton Duc Thang University}, city={Ho Chi Minh City}, country={Vietnam}}

\begin{abstract}
When a code-generating model invents a package name, whoever registers that
name owns every machine that runs the suggestion. The exposure is real but
small, and small quantities are the ones a sloppy pipeline gets wrong. We
measure the rate for one model, and show how easily the measurement
manufactures the number.

We issued \Nprompts{} package-recommendation prompts to \texttt{\Model} under
deterministic decoding, extracted \Npkg{} package names, and checked each
against the live PyPI registry. \Nhall{} do not exist: \Rate\%, 95\% Wilson
interval [\Ratelo\%, \Ratehi\%]. Had the count been zero, the rule-of-three
bound would still have been \Ruleofthree\%, so the measurement is only just
able to distinguish the effect from nothing.

The methodological finding matters more. One answer contained \Ndashchar{}
occurrences of U+2011, a non-breaking hyphen that renders identically to ASCII.
The extraction pattern silently truncated \Ntrunc{} names at that character:
\texttt{pandas-profiling} became \texttt{pandas}. The published total survived
only because every truncated fragment happened to carry the same registry
verdict; two different outcomes would have doubled the reported rate. We give
the corrected rate and the validation steps that would have caught the fault.

Two further measurement choices move the headline as much as the model does.
The unit of analysis is one: the same \Nhall{} events give \Rate\% per
recommended name and \UPromptrate\% per answer, a factor of \URatio, and a
bootstrap that resamples answers rather than names admits zero into the
interval. The characterisation of the suggestion distribution is the other:
its evenness is \SEven\% of maximum and both invented names occur exactly once,
so an attacker who pre-registered both would have intercepted \Rate\% of
recommendations and no more, with no repetition multiplier to apply.
\end{abstract}

\begin{keyword}
Slopsquatting \sep package hallucination \sep software supply chain \sep
LLM code assistants \sep measurement validity \sep Unicode confusables
\end{keyword}

\end{frontmatter}

%% ===========================================================================
\section{Introduction}
\label{sec:intro}

A developer asks an assistant which package to install, receives three
suggestions, and pastes them into a terminal. If one of those names does not
exist, the install fails harmlessly: unless somebody has registered it
first. Registering names that models invent, and waiting for the installs to
arrive, is the attack that has come to be called
\emph{slopsquatting}~\citep{spracklen2025packages}. It inherits its economics
from typosquatting~\citep{vu2020typosquatting, taylor2020spellbound} but
replaces the attacker's guess about human typing errors with something better:
a model that will generate the same wrong name for many users, deterministically,
on demand.

The exposure therefore depends on one empirical quantity: how often a model
names a package that does not exist --- and on how concentrated those inventions
are. Both are measurable. Both are small. This paper measures them for one
model, and spends most of its length on a problem that measuring them made
unavoidable: at rates this low, the error introduced by the measurement pipeline
is the same size as the effect.

\paragraph{What we did.}
We issued \Nprompts{} package-recommendation prompts to \texttt{\Model} through
a commercial inference API under deterministic decoding, in a crossed design of
twelve task stems and ten style suffixes. From the \Nprompts{} answers we
extracted \Npkg{} package names and queried each against the live PyPI JSON
API. \Nhall{} names returned 404.

\paragraph{What we found.}
The hallucination rate is \Rate\% (\Nhall{} of \Npkg{}), 95\% Wilson interval
[\Ratelo\%, \Ratehi\%]. It is a real but small effect, and the interval is wide
in relative terms: the upper bound is more than three times the point estimate.
The recommendations are heavily concentrated (\Ndistinctpkg{} distinct names
across \Npkg{} suggestions, with the ten most frequent accounting for
\Conc\%) which matters, because concentration is what makes the attack worth
an adversary's time.

\paragraph{What went wrong, and why we are reporting it.}
One answer of \Nprompts{} was written with U+2011, the Unicode non-breaking
hyphen. It is visually indistinguishable from the ASCII hyphen in every font we
have looked at. Our extraction pattern used the character class
\texttt{[a-zA-Z0-9\_\-.]}, which does not contain it, so the regular expression
stopped there and returned a prefix:

\begin{center}
\texttt{pandas}\,$\leftarrow$\,\texttt{pandas-profiling}\qquad
\texttt{feature}\,$\leftarrow$\,\texttt{feature-tools}\qquad
\texttt{category}\,$\leftarrow$\,\texttt{category-encoders}
\end{center}

\Ntrunc{} of \Npkg{} extractions were corrupted by one character in one answer.
The study's entire signal is \Nhall{} of \Npkg{}. The two magnitudes are the
same.

The reported total nevertheless came out correct, and it is worth being precise
about why: all \Ncoin{} truncated fragments happen to carry the same PyPI
verdict as the full names they came from. \texttt{pandas} and
\texttt{category} are both real packages, as are \texttt{pandas-profiling} and
\texttt{category-encoders}; \texttt{feature} is absent, as is
\texttt{feature-tools}. The extractor was wrong three times and the arithmetic
was right anyway. Had \texttt{category} not existed on PyPI: and it is not
obvious that it should --- the published rate would have been \CFhi\% rather
than \CFlo\%.

\paragraph{Why this generalises beyond our own bug.}
Non-ASCII look-alike characters are a studied security problem in URLs, in
identifiers and in source code~\citep{gabrilovich2002homograph,
holgers2006homograph, boucher2022badchars, boucher2023trojansource,
unicode-tr36}. What has not been examined, as far as we can tell, is their
effect on the \emph{measurement} pipelines that security papers build around
model outputs. Any study that extracts a structured token from free-form
generated text (package names, CVE identifiers, file paths, function
signatures, URLs) is exposed to exactly this failure, and the smaller the
effect being measured, the larger the relative damage.

\paragraph{Contributions.}
\begin{enumerate}
\item \textbf{A clean measurement.} \Nprompts{} prompts, \Npkg{} extracted
names, every one checked against the live registry;
\Rate\% hallucinated with an explicit interval and an explicit rule-of-three
comparison (Section~\ref{sec:results}).

\item \textbf{An extractor-validity analysis.} We quantify how much of the
reported quantity was decided by the character class in a regular expression,
and give the counterfactual band [\CFlo\%, \CFhi\%] that the uncertainty
implies (Section~\ref{sec:extractor}).

\item \textbf{A concentration analysis.} \Ndistinctpkg{} distinct names, top-ten
share \Conc\%, which is the property an attacker actually exploits
(Section~\ref{sec:concentration}).

\item \textbf{A re-query at a lag.} All \Recheckn{} distinct names were checked
again \Recheckdays{} days later: \Recheckflips{} verdicts changed and both
invented names remain unregistered, so the exposure is open rather than
historical (Section~\ref{sec:recheck}).

\item \textbf{A unit-of-analysis analysis.} The same \Nhall{} events give
\Rate\% per recommended name and \UPromptrate\% per answer, a factor of
\URatio, with non-overlapping practical readings. We estimate the intraclass
correlation across answers ($\mathrm{ICC} = \UIcc$, design effect \UDeff) and
give a cluster bootstrap interval whose lower bound is zero, then state which
unit we report and why (Section~\ref{sec:unit}).

\item \textbf{A validation checklist} for studies that parse model output, and
an explicit statement of what a sub-one-percent rate does not license
(Sections~\ref{sec:checks} and~\ref{sec:threats}).
\end{enumerate}

%% ===========================================================================
\section{Background}
\label{sec:background}

\subsection{Package registries as an attack surface}

Public package registries have been an attack surface for as long as they have
existed~\citep{cappos2008mirror}. The catalogue of realised attacks is
substantial~\citep{ohm2020backstabber}, the taxonomy is
mature~\citep{ladisa2023sok}, and large-scale measurements exist for both
npm~\citep{zimmermann2019npm, zahan2022weaklinks} and
PyPI~\citep{ruohonen2021pypi, alfadel2023python, kaplan2021survey}. Name-based
attacks --- typosquatting and combosquatting --- are among the cheapest, because
they require no compromise of any existing package~\citep{vu2020typosquatting,
duan2021measuring}, and defences based on name-distance
heuristics~\citep{taylor2020spellbound} assume the attacker is modelling human
typing.

\subsection{Slopsquatting}

Slopsquatting changes the generator of the wrong name from a human finger to a
language model. \citet{spracklen2025packages} give the first comprehensive
measurement across many models and prompt sets, and establish the two properties
that make the attack viable: models hallucinate package names at non-trivial
rates, and the hallucinations repeat. Repetition is the crux. A name invented
once is worthless to an attacker; a name invented for many users is an
acquisition target.

Our study is much smaller than theirs in model coverage and is not intended to
compete with it as a survey. It contributes on a different axis: the fidelity of
the measurement instrument at the low rates that current models exhibit.

\subsection{Why models invent package names}

Hallucination in generated text is well characterised in
general~\citep{ji2023hallucination, huang2025hallucination, maynez2020faithfulness}
and specifically as a tendency to produce fluent, plausible tokens in place of
rare factual ones~\citep{lin2022truthfulqa}. Package names sit squarely in the
long tail: \citet{kandpal2023longtail} show that model accuracy on a fact
degrades sharply with how rarely it appears in pre-training, and the PyPI
namespace is overwhelmingly long tail. The pattern we observe is consistent
with that account. Both of our hallucinated names are near-misses of real
packages --- \texttt{feature-tools} for the real \texttt{featuretools}: rather
than inventions from nothing, which is the failure mode one expects when a model
is reconstructing a rarely seen string from its parts.

That near-miss structure interacts badly with registry naming rules. PyPI
normalises names by lowercasing and collapsing runs of \texttt{-},
\texttt{\_} and \texttt{.} to a single \texttt{-}~\citep{pep503, pep508}, so
\texttt{feature\_tools} and \texttt{feature-tools} are the same name, but
\texttt{featuretools} is a different one. A model that inserts a hyphen produces
a genuinely unregistered name, and an attacker can register it.

\subsection{Code assistants and their consumers}

The security literature on code assistants has concentrated on the code they
emit~\citep{chen2021codex, pearce2022asleep, bhatt2023cyberseceval} and on the
developers who accept it~\citep{perry2023insecure, sandoval2023lostatc}, and a
separate line studies deliberate poisoning of the suggestion
model~\citep{schuster2021autocomplete, aghakhani2024trojanpuzzle}. Package
hallucination is different from all three: nobody attacks the model, and the
model is not wrong about security: it is wrong about a string, and the
consequence is a supply-chain compromise. The relevant risk taxonomies place it
under insecure output handling and supply-chain
risk~\citep{owasp2025llm, vassilev2023aml}.

\subsection{The unit of analysis in generated-output studies}

A study that prompts a model $m$ times and extracts $n$ tokens has two
denominators available and they answer different questions. The distinction is
standard in survey sampling, where the tokens are elements and the prompts are
clusters, and the cost of ignoring it is the design effect: the variance of a
proportion estimated from clustered elements exceeds the binomial variance by
$1+(\bar{n}-1)\,\mathrm{ICC}$, where $\mathrm{ICC}$ is the intraclass
correlation of the outcome within clusters and $\bar{n}$ the mean cluster
size~\citep{kish1965survey, killip2004icc, donner2000cluster}. An interval
computed on $n$ elements as though they were independent is therefore too
narrow whenever the outcome correlates within a prompt, which it does whenever
the model's answer is coherent: a model that has decided to recommend a
plausible-sounding but nonexistent library is more likely to do so again in the
same answer.

The literature on model-output measurement does not generally report which unit
it uses, and the two rates are not interchangeable: a per-token rate is the
right quantity for comparing models, and a per-answer rate is the right quantity
for estimating what a developer is exposed to. Section~\ref{sec:unit} reports
both for our sample, along with the design effect and a bootstrap that resamples
prompts rather than tokens~\citep{efron1993bootstrap}, and states which we take
as the headline. We raise the point in the background rather than only in the
threats section because it applies to every study of this shape, including the
ones whose rates ours is compared against.

\subsection{Confusable characters}

The security consequences of characters that render alike are old and
well-documented: homograph domains~\citep{gabrilovich2002homograph,
holgers2006homograph}, imperceptible perturbations that change NLP model
behaviour~\citep{boucher2022badchars}, and bidirectional-override sequences that
make source code read differently to a compiler and to a
reviewer~\citep{boucher2023trojansource}. The Unicode Consortium maintains both
a threat analysis and a set of mechanisms for
mitigation~\citep{unicode-tr36, unicode-uts39}.

All of that literature concerns an adversary choosing a confusable to deceive a
human or a parser. Our case has no adversary. The model emitted U+2011 for
typographic reasons of its own, and our parser was simply not written for it.
The exposure is the same.

%% ===========================================================================
\section{Threat Model}
\label{sec:threat}

Fig.~\ref{fig:threat} states the threat model and the measurement together:
the probability that a recommended name is unregistered, and what an attacker
who registers it gains.

\begin{figure}[t]
\centering
\pendingfig{fig_threat_model}{55mm}
\caption{
Overview. \textbf{(A)} Threat model and measurement: the probability that a
recommended package name is unregistered, \Nhall{} of \Npkg{} names checked
live against the registry. \textbf{(B)} The extraction artifact: one Unicode
character truncated \Ntrunc{} names, a perturbation the size of the signal.
}
\label{fig:threat}
\end{figure}

The adversary can register any unclaimed name on a public registry, at
negligible cost and with no privileged access. They can query the same models
developers use. Their objective is to hold a name that a model will recommend
to somebody who installs it without checking.

The quantity that determines their yield is the probability that a recommended
name is unregistered, multiplied by the concentration of those names: a
recommendation stream that invents a different name every time is not worth
registering against, while one that invents the same name repeatedly is. We
estimate the first quantity directly and characterise the second.

%% ===========================================================================
\section{Method}
\label{sec:method}

\subsection{Prompts}

The prompt set is a crossed design of twelve task stems and ten style suffixes,
giving \Nprompts{} distinct prompts. The stems ask for three package names for a
concrete task --- web scraping, tabular machine learning, async HTTP, HTML
parsing, plotting, PDF parsing, API testing, feature engineering, JSON schema
validation, CLI applications, notebook experiment tracking, data versioning.
The suffixes vary presentation only (``Reply with pip install lines only.'',
``Avoid explanations.'', and so on) and never mention correctness, existence or
verification, so the model is not being cued to self-check.

The system prompt is \emph{``Reply with pip install commands only.''} No prompt
mentions the registry, and no prompt asks the model to avoid inventing names.
The measurement is of ordinary recommendation behaviour, not of behaviour under
a warning.

\subsection{Extraction and registry check}

Each answer is matched against the pattern

\begin{center}
\texttt{pip install\textbackslash s+([a-zA-Z0-9\_\textbackslash-.]+)}
\end{center}

taking up to \Ncapfive{} distinct names per answer in order of appearance. The
cap never binds: the most any answer produces is \Nmaxper{}. The word
\emph{distinct} is load-bearing for the denominator. The corpus contains
\Nocc{} \texttt{pip install} occurrences but \Npkg{} checked names, because one
answer names the same package twice under two extras specifications
(\texttt{httpx[async]} and \texttt{httpx[http2]}) and a second query of the same
name adds no information. \Npkg{}, not \Nocc{}, is therefore the denominator of
every rate in this paper.

Every extracted name is then queried against
\url{https://pypi.org/pypi/{name}/json}.
A 200 response counts as existing, anything else as absent. This is a live
check against the registry as of the run date, not a lookup in a cached index.

Section~\ref{sec:extractor} is about the first of those two steps.

\subsection{Model and decoding}

Table~\ref{tab:setup} records the configuration. One call per prompt, no
retries, deterministic decoding.

\begin{table}[t]
\caption{Run configuration, read from \texttt{results/pilot/run\_meta.json} and
the result files.}
\label{tab:setup}
\centering
\small
\begin{tabular}{l>{\raggedright\arraybackslash}p{82mm}}
\toprule
Parameter & Value \\
\midrule
%%ROWS:tab_setup%%
\bottomrule
\end{tabular}
\end{table}

\subsection{Relationship to an earlier pilot}

An earlier pilot in the same repository reported a hallucination rate of
\Legacyrate\% on \Legacyn{} package names. That figure is not comparable and we
do not treat it as a prior. Its package list was fixed in advance and two of its
\Legacyn{} entries were placeholder strings written by hand
(\texttt{totally-fake-pkg-xyz-2026} and \texttt{nonexistent-slopsquat-abc}), so
the quantity it measured was the fraction of a hand-written list that was
hand-written to be fake, not any property of a model. We report it here only
because Figure~\ref{fig:rate} places the two on one axis, and the two orders of
magnitude between them is itself the argument for extracting names from model
output rather than from a list.

%% ===========================================================================
\section{Results}
\label{sec:results}

\subsection{The rate}

Table~\ref{tab:rate} and Figure~\ref{fig:rate} give the headline numbers.
\Nhall{} of \Npkg{} extracted names are absent from PyPI, a rate of \Rate\%
with a 95\% Wilson interval of [\Ratelo\%, \Ratehi\%]. The two names appear in
\Nhallprompts{} of the \Nprompts{} answers.

\begin{table}[t]
\caption{Hallucination rate under three accountings, with 95\% Wilson
intervals. The last row is the bound the same sample size would have supported
had no hallucinated name been found at all.}
\label{tab:rate}
\centering
\small
\setlength{\tabcolsep}{5pt}
\begin{tabular}{lrrl}
\toprule
Quantity & Count & Rate (\%) & 95\% CI \\
\midrule
%%ROWS:tab_rate%%
\bottomrule
\end{tabular}
\end{table}

\begin{figure}[t]
\centering
\includegraphics[width=\linewidth]{figures/fig_rate.pdf}
\caption{The measured rate on a logarithmic axis, against the earlier
hand-listed pilot and against the rule-of-three bound that \Npkg{} observations
would have supported with a count of zero. The error bar is the 95\% Wilson
interval.}
\label{fig:rate}
\end{figure}

Two observations about the size of this effect. First, the rule-of-three
bound~\citep{hanley1983zero} for \Npkg{} observations is \Ruleofthree\%, which
is inside our confidence interval. A study of this size that had found nothing
would have reported an upper bound not far below our point estimate; the
measurement is only just powered to see the effect it is measuring. Second, the
interval is asymmetric and wide in relative terms, which is the expected
behaviour of a binomial proportion near zero and the reason we use Wilson rather
than the normal approximation~\citep{wilson1927interval, agresti1998approximate,
clopper1934confidence}.

\subsection{Which denominator? The unit of analysis moves the headline}
\label{sec:unit}

The rate above is per extracted name, and that choice is not neutral. The
\Npkg{} names are not \Npkg{} independent draws: they arrive in groups, one
group per answer, and the model chooses a group. \UClusterThree{} of the
\UNprompt{} answers name three packages, \UClusterTwo{} names two and
\UClusterOne{} name one, a mean of \UNbar{} names per answer. A rate computed
over names therefore treats a clustered sample as a simple random one, and a
rate computed over answers asks a different question: not ``what fraction of
recommended names do not exist'' but ``what fraction of answers contains at
least one name that does not exist''. The second is the question a developer
faces, since a developer consumes answers rather than names.

Table~\ref{tab:unit} gives all three accountings. The mention-level rate is
\Rate\% with a Wilson interval of [\UMwLo\%, \UMwHi\%]. The answer-level rate is
\UPromptrate\% --- \UPrompthit{} of \UNprompt{} answers: with an interval of
[\UPwLo\%, \UPwHi\%]. The point estimates differ by a factor of \URatio, and the
answer-level upper bound is nearly three times the mention-level one. Neither is
wrong. They are answers to different questions, and a paper that reports one
without naming the unit invites the reader to apply it to the other.

\paragraph{Clustering does not inflate the variance here, and we checked rather
than assumed.} The usual consequence of cluster sampling is that the effective
sample size falls below the nominal one, so the Wilson interval computed on
\Npkg{} would be too narrow. We estimated the one-way intraclass correlation of
the hallucination indicator across answers and obtained
$\mathrm{ICC} = \UIcc$, giving a design effect of
$1 + (\bar{n}-1)\,\mathrm{ICC} = \UDeff$ and an effective sample size of
\UEss{} against a nominal \Npkg{}. The ICC is slightly negative and the design
effect is therefore slightly below one, which means the mention-level interval
is if anything marginally conservative rather than optimistic.

The reason is mechanical and worth stating so that nobody reads the result as a
general licence to ignore clustering. With \Nhall{} events in the whole sample,
the only arrangements available are both events in one answer or one event in
each of two answers. The data realise the second, which is the maximally
dispersed arrangement and the one that minimises within-answer correlation. Had
both hallucinated names landed in the same answer the ICC would have been
positive and the design effect above one. So the finding is that this particular
sample happens not to be clustered in the relevant sense, not that answer-level
clustering is negligible in this setting. At \Nhall{} events the quantity is
barely estimable and we report it as a check that passed rather than as a
measurement.

\paragraph{A resampling interval that respects the clustering.} A bootstrap that
resamples answers rather than names gives a further reading. Over \UBoot{}
resamples the 95\% percentile interval on the mention-level rate is
[\UCbLo\%, \UCbHi\%]. Its lower bound is exactly zero, and that is informative
rather than a defect: because only \UPrompthit{} of \UNprompt{} answers contain
a hallucinated name, a resample of \UNprompt{} answers drawn with replacement
misses both of them often enough to put zero inside the interval. The honest
summary of a two-event sample is that it is consistent with a rate of zero once
the sampling unit is respected, and any bound we quote should be read as an
upper bound rather than as a located estimate.

\paragraph{Which one we report, and why.} We keep the mention-level rate as the
headline, because it is the quantity the title names and the quantity comparable
to the rates the surrounding literature reports, and we state the answer-level
rate beside it throughout. Where the distinction matters for the conclusion we
use the answer-level number, since exposure accrues per answer acted on. What we
do not do is pick whichever unit gives the more striking figure, which in this
data would be the answer level and would triple the headline.

\begin{table}[t]
\caption{The same \Nhall{} hallucinated names under three accountings. Rows
differ only in what counts as one observation. The clustered row resamples
answers rather than names over \UBoot{} bootstrap replicates.}
\label{tab:unit}
\centering
\small
\setlength{\tabcolsep}{5pt}
\begin{tabular}{@{}lrrrl@{}}
\toprule
Unit & Events & $n$ & Rate (\%) & 95\% CI \\
\midrule
%%ROWS:tab_unit%%
\bottomrule
\end{tabular}
\end{table}

\subsection{Where in the answer the invented names appeared}
\label{sec:pos}

Table~\ref{tab:pos} records the position of each hallucinated name within its
answer's list. Both fall at position \UPosBoth{} of three.

We report this and draw nothing from it. With two events, any positional pattern
has a probability under the null of roughly one in nine of arising by chance, so
the observation is not evidence that models are more likely to invent a name in
the middle of a list. We include the column because a reader checking our
extraction will want to know where in the output the two names sat, and because
if a larger replication does find a positional effect, this row is the datum it
would be compared against. Treating it as a finding now would be exactly the
over-reading this paper is about.

\begin{table}[t]
\caption{Position of each hallucinated name within its answer. Recorded for
replication, not interpreted: \Nhall{} events support no positional claim.}
\label{tab:pos}
\centering
\small
\setlength{\tabcolsep}{6pt}
\begin{tabular}{@{}llc@{}}
\toprule
Prompt & Name & Position \\
\midrule
%%ROWS:tab_pos%%
\bottomrule
\end{tabular}
\end{table}

\subsection{The two hallucinated names}

Table~\ref{tab:positives} lists them with the line of model output that
produced each. Neither is an invention from nothing. \texttt{feature-tools} is a
hyphenated variant of the real package \texttt{featuretools}; under PyPI's
normalisation rules~\citep{pep503} the two are distinct names, and the
hyphenated one is unregistered and therefore claimable. \texttt{cleverscope} is
a plausible-sounding name in a list otherwise composed of real tools
(\texttt{dvc}, \texttt{gitpython}).

\begin{table}[t]
\caption{The \Nhall{} extracted names absent from PyPI, with the model output
line that produced each. The bracketed codepoint marks the non-breaking hyphen
discussed in Section~\ref{sec:extractor}.}
\label{tab:positives}
\centering
\small
\setlength{\tabcolsep}{5pt}
\begin{tabular}{llp{58mm}}
\toprule
Name & On PyPI & Model output line \\
\midrule
%%ROWS:tab_positives%%
\bottomrule
\end{tabular}
\end{table}

That both are near-misses rather than fabrications is the practically important
detail. A defence that checks whether a suggested name exists will catch both. A
developer who reads the suggestion will catch neither, because both look exactly
like packages they have heard of.

\subsection{Concentration of recommendations}
\label{sec:concentration}

Across \Npkg{} suggestions the model named \Ndistinctpkg{} distinct packages.
The ten most frequent account for \Conc\% of all suggestions.
Table~\ref{tab:top} and Figure~\ref{fig:top} give the distribution.

\begin{table}[t]
\caption{The twelve most-suggested package names across \Nprompts{} prompts.}
\label{tab:top}
\centering
\small
\begin{tabular}{lrrl}
\toprule
Package & Times suggested & Share (\%) & On PyPI \\
\midrule
%%ROWS:tab_top%%
\bottomrule
\end{tabular}
\end{table}

\begin{figure}[t]
\centering
\includegraphics[width=\linewidth]{figures/fig_top.pdf}
\caption{Suggestion frequency for the fifteen most-recommended names, coloured
by whether the name exists on PyPI. The distribution is dominated by a small
set of well-known packages.}
\label{fig:top}
\end{figure}

\paragraph{``Concentrated'' overstates it.} The top-ten share is the number
usually quoted, and taken alone it suggests a short head and a long tail. The
full coverage curve (Table~\ref{tab:cover}) says something different. Half the
suggestions require the top \SHalfAt{} names and ninety per cent require the top
\SNinetyAt{}, out of \Ndistinctpkg{} distinct names in all. Shannon entropy over
the suggestion distribution is \SShannon{} nat against a maximum of
$\ln \Ndistinctpkg = \SShannonMax$, an evenness of \SEven\%, and the
corresponding effective number of names is $\exp(H) = \SEffNames$. A
distribution in which \SEffNames{} of \Ndistinctpkg{} names are effectively in
play is close to uniform, not concentrated. We correct the framing here because
the attacker-yield argument below depends on which of the two descriptions is
accurate.

\paragraph{What an attacker would actually collect.} Each mention is
\SPerName\% of the total, so a registered name intercepts that share of
recommendations per occurrence. Both hallucinated names occur exactly once, so
an attacker who had pre-registered both would have intercepted
\SHallMentions{} of \Npkg{} recommendations: the same \Rate\% that is the
paper's headline, which is the sense in which the headline \emph{is} the
attacker's yield on this sample. There is no repetition multiplier to apply:
\SOnes{} of the \Ndistinctpkg{} distinct names (\SOnesPctNames\% of names,
\SOnesPctMentions\% of mentions) occur exactly once, and both invented names
fall in that group.

This is the opposite of the dynamic that makes slopsquatting attractive in the
reports that motivated this study, where the same invented name recurs across
prompts and one registration harvests many installs. On this model and this
prompt distribution the invented names do not recur, so the attack degenerates
into registering many names for one interception each. We state this as a
property of our sample rather than of the technique: a model with a stronger
prior toward a particular nonexistent name, or a prompt distribution narrower
than ours, would produce the recurrence we do not observe, and our \Nprompts{}
prompts cannot rule that out.

\begin{table}[t]
\caption{Cumulative share of the \Npkg{} suggestions covered by the $k$ most
frequent names. The last row is all \Ndistinctpkg{} distinct names.}
\label{tab:cover}
\centering
\small
\setlength{\tabcolsep}{6pt}
\begin{tabular}{@{}rrr@{}}
\toprule
$k$ & Suggestions covered & Share (\%) \\
\midrule
%%ROWS:tab_cover%%
\bottomrule
\end{tabular}
\end{table}

\subsection{The window is still open}
\label{sec:recheck}

A hallucinated name is only a vulnerability while it remains unregistered, so
the registry verdict is a measurement with a shelf life. We re-queried every
one of the \Recheckn{} distinct names \Recheckdays{} days after the run, on
\Recheckdate.

Nothing moved. \Recheckflips{} of the \Recheckn{} verdicts changed, and both
invented names --- \texttt{feature\-tools} and \texttt{cleverscope}: are
still absent from PyPI. Neither has been claimed in the interval, by a defender
reserving it or by anyone else.

That is one observation at one lag and we draw no rate from it. What it
establishes is narrower and still worth stating: the exposure this paper
measures is not hypothetical and has not been closed. A reader reproducing this
work should re-run the check rather than reuse our verdicts, and a defender who
wants these two names neutralised can register them today.

\subsection{Determinism}

Decoding is deterministic and the ten style suffixes are semantically empty, so
the \Nprompts{} prompts produced only \Ndistinctout{} distinct answers. This
limits how much independent evidence the design carries: repeated
near-identical prompts under $T=0$ do not contribute independent observations.
We report the interval on \Npkg{} extracted names because that is the unit the
registry check applies to, but a reader should treat the effective number of
independent trials as closer to \Ndistinctout{} than to \Nprompts{}, and the
intervals in Table~\ref{tab:rate} as correspondingly optimistic.

%% ===========================================================================
\section{The Extractor Is Part of the Measurement}
\label{sec:extractor}

\subsection{What happened}

One answer of \Nprompts{} contained \Ndashchar{} occurrences of U+2011, the
Unicode non-breaking hyphen. The extraction pattern's character class contains
U+002D and not U+2011, so each match terminated at the first such character and
returned a prefix. Table~\ref{tab:extractor} shows the three affected names, the
fragments the extractor returned, and the registry verdict for each.

\begin{table}[t]
\caption{The \Ntrunc{} extractions corrupted by a single non-breaking hyphen in
one model answer, with the live PyPI verdict for the full name and for the
truncated fragment.}
\label{tab:extractor}
\centering
\small
\setlength{\tabcolsep}{5pt}
\begin{tabular}{llccl}
\toprule
Name the model wrote & Extracted & On PyPI & Fragment on PyPI & Verdict \\
\midrule
%%ROWS:tab_extractor%%
\bottomrule
\end{tabular}
\end{table}

\subsection{Why the total was right anyway}

All \Ncoin{} fragments carry the same registry verdict as the names they came
from. That is a coincidence of the PyPI namespace, not a property of the
extractor: \texttt{pandas} and \texttt{category} both happen to be registered
names, and \texttt{feature} happens not to be. Under a namespace in which
\texttt{category} were unclaimed: a single registration away --- the same
run would have reported \Nhallcorr{}$+1$ hallucinated names.

\subsection{The counterfactual band}

Figure~\ref{fig:extractor} shows the rate under three accountings: the ASCII
extractor as run, the Unicode-aware extractor after repair, and the worst case
in which the coincidences had gone the other way. The band is
[\CFlo\%, \CFhi\%]: a factor of two, produced by one character in one of
\Nprompts{} answers.

\begin{figure}[t]
\centering
\includegraphics[width=\linewidth]{figures/fig_extractor.pdf}
\caption{The reported rate under three extraction accountings, with 95\% Wilson
intervals. The gap between the first two bars is zero only by coincidence; the
third bar is what the same data would have yielded had the coincidence gone the
other way.}
\label{fig:extractor}
\end{figure}

\subsection{Where the same character class was right}

The extras case of Section~\ref{sec:method} is the same mechanism with the
opposite outcome, and it is worth stating because it bounds the complaint. The
character class excludes \texttt{[} as well as U+2011, so
\texttt{pip install httpx[async]} also terminates early and returns a prefix.
Here the prefix is the correct answer: \texttt{httpx} is the distribution name
and \texttt{[async]} selects optional dependencies of it, so the truncation
recovers exactly what should be queried. \Nextras{} of \Nprompts{} answers use
that syntax.

The two cases are indistinguishable from inside the extractor. Both are
``the match stopped at a character the class does not contain''; one is a bug
and one is the intended behaviour, and which it is depends on what the character
means in the grammar of the token, not on anything the pattern can see. An
extractor written as a character class therefore cannot be validated by reading
it. It has to be validated against its output, which is the point of
Section~\ref{sec:checks}.

\subsection{The general form of the problem}

The failure is not specific to hyphens, to PyPI, or to this model. It has three
ingredients, and any study with all three is exposed:

\begin{enumerate}
\item A structured token is extracted from free-form generated text by a pattern
written over an assumed character repertoire.
\item The generator is free to emit characters outside that repertoire that
render identically to characters inside it.
\item The quantity being estimated is small, so a handful of corrupted
extractions is comparable to the whole effect.
\end{enumerate}

Package names satisfy all three, and so do CVE identifiers, file paths, function
signatures, domain names and version specifiers. The Unicode Consortium's
confusables data~\citep{unicode-uts39} enumerates the characters that satisfy
the second condition; there are thousands of them.

%% ===========================================================================
\section{Validation Steps That Would Have Caught This}
\label{sec:checks}

Each of the following is a few lines of code and runs in the time it takes to
read the result file.

\paragraph{Check 1: assert the output is ASCII, or normalise before matching.}
Scan every generated answer for codepoints outside the repertoire the extractor
assumes and fail loudly. Ours would have failed on one answer out of
\Nprompts{}, at which point the truncation is obvious. Where failing is too
strict, normalise the confusable classes first: mapping the dash block to
U+002D is a one-line substitution.

\paragraph{Check 2: extract twice and diff.}
Run a permissive extractor alongside the strict one and assert that the two
agree. Disagreements are exactly the cases where the character class is doing
work you did not intend. This is how we found the fault, after the fact.

\paragraph{Check 3: eyeball every positive.}
At a rate below one percent the positives are few enough to read. Both of ours
are one line of output each. Reading them would have shown
\texttt{feature} sitting next to a visible hyphen and a suffix that the table
did not contain.

\paragraph{Check 4: report the count, not only the rate.}
``\Rate\%'' hides that the numerator is \Nhall{}. A reader who sees
``\Nhall{} of \Npkg{}'' immediately knows the estimate is fragile in a way that
a percentage to two decimal places conceals.

\paragraph{Check 5: state the extractor in the paper.}
The regular expression is a parameter of the experiment, not an implementation
detail. Ours appears in Section~\ref{sec:method}, and the substantive finding of
this paper is about it.

%% ===========================================================================
\section{Discussion}
\label{sec:discussion}

\subsection{What a rate near half a percent means}

It means that in a stream of package recommendations from this model on this
prompt distribution, roughly one name in two hundred is claimable. Whether that
constitutes a serious risk depends on volume rather than on rate: an assistant
serving a million suggestions a day emits thousands of claimable names a day at
this rate. The rate is small; the exposure is not necessarily.

It also means that any comparison between models, or any claim that a
mitigation reduced the rate, needs a sample size that our study does not have.
With \Npkg{} observations the interval spans a factor of thirteen. Distinguishing
$0.5\%$ from $0.25\%$ at conventional power requires observations in the tens of
thousands.

\subsection{Near-misses, not inventions}

Both hallucinated names are perturbations of real packages. This supports the
long-tail account of hallucination~\citep{kandpal2023longtail} over an account
in which the model fabricates freely, and it has a defensive consequence: the
namespace an attacker should pre-register is not arbitrary strings but the
hyphenation, pluralisation and separator variants of popular packages. That set
is enumerable and small, which suggests a registry-side defence: reserving
normalisation-adjacent variants of high-download packages --- that does not
depend on detecting anything about the model.

\subsection{Measurement fragility as a finding}
The literature on this attack is young, and the rates being reported are small.
A field in which the headline quantity is a few tenths of a percent needs
instrument validation as a routine part of its method sections, in the way that
detection thresholds and false-positive rates are routine in intrusion
detection. We would rather this paper be read as an argument for that practice
than as a rate for one model.

%% ===========================================================================
\section{Threats to Validity}
\label{sec:threats}

\paragraph{One model, one prompt distribution, one day.} All \Nprompts{} calls
used \texttt{\Model} on a single date. Package existence is time-varying: a name
absent today may be registered tomorrow, including by an attacker who read this
paper. The registry verdicts are a snapshot.

\paragraph{Small numerator.} The estimate rests on \Nhall{} events. Every
interval in this paper should be read with that in mind, and no comparison
against another model or another study is supported by this sample. The
cluster bootstrap of Section~\ref{sec:unit} makes the point quantitatively:
resampling answers puts \UCbLo\% inside the 95\% interval, so the data are
consistent with the model never inventing a package name and our having drawn
two unlucky answers. We regard the upper bound as the defensible output of this
measurement and the point estimate as provisional.

\paragraph{The design effect we report is barely estimable.} $\mathrm{ICC} =
\UIcc$ on \Nhall{} events is a number computed on almost no information, and we
present it as a check rather than an estimate. Its negative sign follows
mechanically from the two events falling in different answers, which is the
only dispersed arrangement available; it is not evidence that hallucinations are
negatively correlated within an answer. A study designed to estimate the ICC
would need enough events to populate both arrangements, which at this rate means
on the order of a thousand prompts.

\paragraph{Deterministic decoding.} At $T=0$ the \Nprompts{} prompts yielded
\Ndistinctout{} distinct answers, so the effective sample is smaller than the
nominal one. A temperature sweep would give a fuller picture of the
hallucination distribution and is not attempted here.

\paragraph{Extraction ceiling.} The extractor takes at most five names per
answer and requires the literal string \texttt{pip install}. Recommendations
phrased differently are invisible to it. This biases the sample toward
well-formatted answers.

\paragraph{Existence is not the only harm.} A name that exists may still be
malicious, abandoned, or wrong for the task. We measure only whether it
resolves.

\paragraph{The corrected numbers are ours.} The repair in
Section~\ref{sec:extractor} re-queried PyPI for the affected names and froze
those verdicts into the analysis code so the paper reproduces offline. A reader
re-running the study will query a registry that has since changed.

%% ===========================================================================
\section{Conclusion}

We measured how often one code-generating model names a Python package that
does not exist: \Nhall{} of \Npkg{} extracted names, \Rate\%, 95\% CI
[\Ratelo\%, \Ratehi\%]. Both hallucinations are near-misses of real packages
rather than inventions, which points at a cheap registry-side defence.

The finding we would rather be remembered for is smaller and less comfortable.
One invisible character in one of \Nprompts{} answers corrupted \Ntrunc{} of
\Npkg{} extractions in a study whose whole signal is \Nhall{} of \Npkg{}. The
published total was correct by coincidence, and the counterfactual band the
fault admits is [\CFlo\%, \CFhi\%]. When the quantity being estimated is this
small, the parser is not infrastructure. It is the experiment.

%% ===========================================================================
\appendix

\section{Reproduction}
\label{app:repro}

\subsection{What the pipeline reads and writes}

Table~\ref{tab:inputs} lists them.

\begin{table}[h]
\caption{Input files. Sizes are of the released artefacts; the raw answers are
the only file a reader needs in order to re-derive everything else.}
\label{tab:inputs}
\centering
\small
\setlength{\tabcolsep}{4pt}
\begin{tabular}{@{}>{\raggedright\arraybackslash}p{0.44\linewidth}>{\raggedright\arraybackslash}p{0.48\linewidth}@{}}
\toprule
Path under \texttt{results/} & Contents \\
\midrule
\texttt{pilot/raw\_model\_outputs.jsonl} &
  the \Nprompts{} verbatim answers, one JSON object each \\
\texttt{pilot/package\_validation.csv} &
  one row per extracted name: prompt, name, registry verdict \\
\texttt{pilot/metrics.csv} &
  the pilot's own summary line \\
\texttt{pilot/run\_meta.json} &
  provider, model, decoding, call count, wall-clock \\
\texttt{recheck/pypi\_recheck.json} &
  live status of all \Recheckn{} distinct names, re-queried \Recheckdays{} days
  later \\
\bottomrule
\end{tabular}
\end{table}

\subsection{The two extractors, verbatim}

The paper's central methodological point is that a character class decided part
of the reported quantity, so we give both patterns rather than describe them.
The pilot used
\begin{center}
\verb|pip install\s+([a-zA-Z0-9_\-\.]+)|
\end{center}
and the repaired analysis uses
\begin{center}
\verb|pip install\s+([^\s=<>[],;'"]+)|
\end{center}
The first stops at any character outside its class, so a U+2011 non-breaking
hyphen truncates the name; the second stops only at characters a shell would
treat as a token boundary. Both are applied case-insensitively to the raw
answers, and \texttt{analysis.py} runs both on every answer so the difference
between them is computed rather than asserted.

\subsection{Which numbers are measured and which are derived}

Measured: the \Nprompts{} answers, the registry verdict for each name, the
run metadata. Derived: every rate, every interval, the re-extraction under the
widened pattern, the concentration figures, the intraclass correlation and
design effect, and the cluster bootstrap. The derivation is one code path ---
\texttt{analysis.py}: imported by both \texttt{make\_tables.py} and
\texttt{make\_figures.py}, so a table and a figure showing the same quantity
cannot disagree.

Two constants in the analysis are choices rather than measurements and we name
them here: the \Ncapfive{}-name cap the extractor applies per answer, and the
\UBoot{} replicates of the cluster bootstrap. Both are stated at the point of
use as well.

\subsection{Offline reproducibility, and its cost}

The registry verdicts for the names the two extractors disagree on are frozen
into \texttt{analysis.py} as a literal dictionary, queried on 5 September 2026.
This is what lets the paper rebuild without network access, and it is also a
limitation: a reader who re-queries PyPI today may get different verdicts, and
the frozen values are the ones the paper's arithmetic uses. The re-query
recorded in \texttt{recheck/pypi\_recheck.json} is the independent check on that
freezing, and it found \Recheckflips{} verdict changes across all \Recheckn{}
distinct names.

\section*{CRediT authorship contribution statement}

\textbf{Son X. Ha:} Conceptualization, Formal analysis, Investigation, Writing -- original draft, Writing -- review \& editing.
\textbf{Phat T. Tran-Truong:} Methodology, Validation, Formal analysis.
\textbf{Xuan-Bach Le:} Software, Validation, Investigation.
\textbf{Trung Phan Hoang Tuan:} Supervision, Resources, Writing -- review \& editing.
\textbf{Tung Q. Nguyen:} Software, Data curation, Visualization, Investigation.
\textbf{Tran Gia Huy:} Investigation, Data curation.
\textbf{Nhan Nguyen Tran Kha:} Investigation, Visualization.
\textbf{Thuc Nguyen Minh:} Investigation, Data curation.
\textbf{Nguyen Ngoc Phien:} Supervision, Project administration, Writing -- review \& editing.

\section*{Declaration of Competing Interest}

The authors declare that they have no known competing financial interests or
personal relationships that could have appeared to influence the work reported
in this paper.

\section*{Funding}

This research received no specific grant from any funding agency in the
public, commercial or not-for-profit sectors.

\section*{Data Availability}

The \Nprompts{} verbatim model answers, the \Npkg{} extracted names with their
live registry verdicts, the \Recheckdays-day re-query of all \Recheckn{}
distinct names, and the run metadata recording model and decoding parameters are
released with the paper, together with the analysis scripts that produce every
number and figure in it. Registry verdicts are timestamped, because they are
observations of a mutable namespace rather than constants.

The artefact (result files, generator scripts and \LaTeX{} sources) is not
publicly posted at the time of submission. We will deposit it in a public
repository under a persistent identifier upon acceptance, and we can supply it
to the editorial office at any point during review.

\bibliographystyle{elsarticle-num}
\bibliography{refs}

\end{document}
